\section{Gauges and projection measures}\label{sec:gauges}

We first associate a measure to an arbitrary convex body.  A volume
variation will give both the lower bound satisfied by this measure and
its relation to the projection body.  No symmetry of the original body
is needed.

\subsection{Differentiating a gauge}

For a convex body $L$ containing the origin in its interior, its
\emph{gauge} is
\[
 g_L(x)=\inf\{a>0:x\in aL\}.
\]
It is convex and positively homogeneous, and $L=\{x:g_L(x)\le1\}$.
When $L$ is origin-symmetric, its gauge is a norm.  A continuous function
$\phi$ on $\Sph$, when evaluated at a vector, will mean its positively
homogeneous extension
\[
 \phi(x)=|x|\phi(x/|x|)\quad(x\ne0),\qquad \phi(0)=0.
\]
This convention does not impose evenness on $\phi$.

Varying the supporting constraints in this way is a Wulff variation.
The almost-everywhere radial differentiation below is the gauge form
of the Aleksandrov variation used in dual Brunn--Minkowski theory;
compare \cite[Lemmas~4.1--4.3 and Corollary~4.9]{HuangLutwakYangZhang2016}. We give the
argument directly because the later minimization uses arbitrary
continuous constraints of both signs, without boundary regularity.

\begin{lemma}[Gauge variation]\label{lem:gauge-variation}
Let $s:\Sph\to(0,\infty)$ and $\phi:\Sph\to\R$ be continuous.  For
sufficiently small $|t|$, define
\[
 W_t=\bigcap_{v\in\Sph}\{x\in\R^n:x\cdot v\le s(v)+t\phi(v)\},
 \qquad g_t=g_{W_t}.
\]
Then $W_t$ is a convex body containing the origin in its interior, and
\begin{equation}\label{eq:gauge-formulas}
 g_t(x)=\max_{v\in\Sph}\frac{x\cdot v}{s(v)+t\phi(v)},
 \qquad
 \frac{|W_t|}{\kappa_n}=\E g_t(u)^{-n}.
\end{equation}
For almost every $u\in\Sph$, the maximizing vector $v/s(v)$ at $t=0$
is unique and equals the Euclidean gradient $\nabla g_0(u)$.  At these
directions,
\begin{equation}\label{eq:gauge-derivative}
 \left.\frac{d}{dt}\right|_{t=0}g_t(u)
 =-g_0(u)\phi\bigl(\nabla g_0(u)\bigr).
\end{equation}
On the sphere the gauges are uniformly bounded above and away from
zero for small $|t|$, and the difference quotients
$(g_t-g_0)/t$ are uniformly bounded.  In particular,
\begin{equation}\label{eq:constraint-volume-derivative}
 \left.\frac{d}{dt}\right|_{t=0}|W_t|
 =n\kappa_n\E\!\left[
   g_0(u)^{-n}\phi\bigl(\nabla g_0(u)\bigr)\right].
\end{equation}
For $s=h_L$, one has $W_0=L$.
\end{lemma}

\begin{proof}
Choose $0<m\le s(v)\le M$.  For small $|t|$ the constraints
$s+t\phi$ lie between $m/2$ and some fixed finite constant $M'$.
Consequently
\[
 (m/2)B_2^n\subset W_t\subset M'B_2^n.
\]
The defining inequalities show that $x\in aW_t$ precisely when
$x\cdot v\le a(s(v)+t\phi(v))$ for every $v$.  Taking the least
such $a$ gives the maximum formula in \eqref{eq:gauge-formulas}.
The radial function of $W_t$ is $1/g_t(u)$, so polar integration,
with spherical area $n\kappa_n$, gives the volume formula.  If
$s=h_L$, separation identifies the defining intersection with $L$.

For completeness, the maximum formula also gives the differentiability
needed here without any regularity assumption on the boundary of
$W_0$.  The function $g_0$ is convex and Lipschitz: it is the maximum
of linear functions whose coefficient vectors $v/s(v)$ form a compact
set.  On every coordinate line, a finite convex function has a
derivative except on a set of one-dimensional measure zero.  Slicing
therefore shows that all coordinate partial derivatives of $g_0$
exist at almost every point of $\R^n$.  If a vector $w=v/s(v)$ attains
the maximum at such a point $x$, then
\[
 g_0(x+he_i)\ge g_0(x)+h w_i\qquad(h\in\R).
\]
The two signs of $h$ imply $w_i=\partial_i g_0(x)$ for every $i$.
Thus the maximizing vector is unique.

Uniqueness implies differentiability.  Indeed, if $w_z$ maximizes at
$x+z$ and $w$ maximizes at $x$, compactness and continuity imply
$w_z\to w$ as $z\to0$.  The maximum formula gives
\[
 0\le g_0(x+z)-g_0(x)-w\cdot z
 \le (w_z-w)\cdot z=o(|z|).
\]
The nondifferentiability set away from the origin is invariant under
positive dilations, by homogeneity.  Since it has Lebesgue measure
zero, polar integration shows that its intersection with $\Sph$ has
$\sigma$-measure zero.

Fix a differentiability direction $u$, and let $v_0$ maximize at
$t=0$.  Uniqueness of $v_0/s(v_0)$ implies uniqueness of $v_0$
itself, since $s$ is positive.  Every maximizing direction $v_t$ for
$g_t(u)$ tends to $v_0$ as $t\to0$.  The derivatives with respect
to $t$ of the functions
\[
 (u,v,t)\longmapsto\frac{u\cdot v}{s(v)+t\phi(v)}
\]
are continuous on the relevant compact set.  Comparing the maxima at
$t$ and $0$, using $v_t$ and $v_0$, therefore gives
\[
 \left.\frac{d}{dt}\right|_{0}g_t(u)
 =-\frac{(u\cdot v_0)\phi(v_0)}{s(v_0)^2}
 =-g_0(u)\phi\!\left(\frac{v_0}{s(v_0)}\right).
\]
This is \eqref{eq:gauge-derivative}.  Notice that the argument uses
neither convexity nor evenness of the function $s+t\phi$.

The same uniform bound on the derivatives of the displayed quotients
bounds $(g_t-g_0)/t$ uniformly in $u$.  The ball inclusions bound
$g_t$ above and away from zero on $\Sph$.  Dominated convergence
now differentiates the volume formula and proves
\eqref{eq:constraint-volume-derivative}.
\end{proof}

\subsection{A volume inequality with its equality case}

We use the classical Brunn--Minkowski inequality, including its equality
case, in the following form; see \cite[Sections~3 and~5]{Gardner2002}.
The first-variation statement and its equality consequence are then
proved in the gauge coordinates needed for our measures.

\begin{theorem}[Brunn--Minkowski inequality]\label{thm:brunn-minkowski}
For every integer $d\ge1$ and convex bodies $A,B\subset\R^d$, write $|\cdot|$ for
$d$-dimensional volume. For $0\le t\le1$,
\begin{equation}\label{eq:brunn-minkowski}
 |(1-t)A+tB|^{1/d}
 \ge (1-t)|A|^{1/d}+t|B|^{1/d}.
\end{equation}
If $0<t<1$, equality holds if and only if $B=a+cA$ for some
$a\in\R^d$ and $c>0$.
\end{theorem}



\begin{corollary}[First variation of volume]\label{cor:volume-first-variation}
Let $K,M\subset\R^n$ be convex bodies, with the origin interior to
$K$.  Then
\begin{align}
 \left.\frac{d}{dt}\right|_{0+}|K+tM|
 &=n\kappa_n\E\!\left[g_K(u)^{-n}
                 h_M\bigl(\nabla g_K(u)\bigr)\right]
                 \label{eq:minkowski-volume-derivative}\\
 &\ge n|K|^{(n-1)/n}|M|^{1/n}.\nonumber
\end{align}
Equality in the inequality holds if and only if $K$ and $M$ are
homothetic up to translation.  The derivative identity remains valid
when $M$ is an arbitrary nonempty compact convex set.
\end{corollary}

\begin{proof}
Support functions add under Minkowski addition, so
Lemma~\ref{lem:gauge-variation}, with $s=h_K$ and $\phi=h_M$,
gives the derivative identity.  This applies equally to a compact
convex set $M$ of lower dimension.  For full-dimensional $M$,
Theorem~\ref{thm:brunn-minkowski} and homogeneity give
\[
 F(t):=|K+tM|^{1/n}\ge |K|^{1/n}+t|M|^{1/n}
 \qquad(t\ge0).
\]
In addition, $F$ is concave, by applying that theorem to $K+sM$
and $K+tM$.  Its right derivative at zero therefore satisfies
$F'(0+)\ge |M|^{1/n}$, which is the asserted lower bound.

If equality holds in this derivative bound, concavity gives
\[
 F(t)\le F(0)+tF'(0+)=|K|^{1/n}+t|M|^{1/n}.
\]
Thus the preceding volume inequality is an equality for every $t>0$.
The equality case of Theorem~\ref{thm:brunn-minkowski}, applied for
example to $K+M$, forces homothety.  Conversely, if $M=a+cK$
with $c>0$, then $K+tM=ta+(1+tc)K$, which verifies equality in
the derivative bound.
\end{proof}

\subsection{The measure associated with a convex body}

The construction below is a gauge-coordinate form of the classical
surface-area and mixed-volume description of projection bodies;
see \cite[Section~6]{Gardner2002} and \cite[Section~1]{Saroglou2015}.
The local segment calculation will fix its normalization and retain
the equality information from the first variation.

For a finite positive Borel measure $\eta$ on $\R^n$ with bounded
support, write
\[
 \tau_\eta(M)=\int h_M(x)\,d\eta(x).
\]
\begin{definition}\label{def:admissible-measure}
We call $\eta$ \emph{admissible} if its support spans $\R^n$ and
\begin{equation}\label{eq:functional-condition}
 \tau_\eta(M)\ge\left(\frac{|M|}{\kappa_n}\right)^{1/n}
 \quad\text{for every origin-symmetric convex body }M\subset\R^n.
\end{equation}
The measure itself is not required to be symmetric.
\end{definition}

\begin{proposition}[Projection measure]\label{prop:projection-measure}
Translate and dilate a convex body $K$ so that $0\in\operatorname{int}K$
and $|K|=\kappa_n$.  Define $\eta_K$ by
\begin{equation}\label{eq:eta-definition}
 \int \psi(x)\,d\eta_K(x)
 =\E\!\left[g_K(u)^{-n}\psi\bigl(\nabla g_K(u)\bigr)\right]
\end{equation}
for nonnegative Borel functions $\psi$; the choice of the gradient on
its null exceptional set is immaterial.  Then $\eta_K$ is an
admissible probability measure of bounded support.  Equality in
\eqref{eq:functional-condition} for $\eta_K$ and a symmetric body
$M$ holds if and only if $K$ is a translate of a positive dilate of
$M$.  Moreover,
\begin{equation}\label{eq:projection-measure}
 h_{\Pi K}(y)=\frac{n\kappa_n}{2}
                  \int |x\cdot y|\,d\eta_K(x)
 \qquad(y\in\R^n),
\end{equation}
and
\begin{equation}\label{eq:eta-normalization}
 \int h_K(x)\,d\eta_K(x)=1.
\end{equation}
For the unit ball, $\eta_{B_2^n}=\sigma$, and consequently
\begin{equation}\label{eq:ball-constant}
 \frac{n\kappa_n b_n}{2}=\kappa_{n-1}.
\end{equation}
\end{proposition}

\begin{proof}
By \eqref{eq:gauge-formulas}, $\E g_K^{-n}=1$, so
\eqref{eq:eta-definition} defines a probability measure.  The active
vectors $v/h_K(v)$ lie in a bounded set because $h_K$ is positive
on the sphere.  Hence $\eta_K$ has bounded support.  Applying
Corollary~\ref{cor:volume-first-variation} gives, for every symmetric
convex body $M$,
\[
 \tau_{\eta_K}(M)
 =\frac1{n\kappa_n}
        \left.\frac{d}{dt}\right|_{0+}|K+tM|
 \ge\left(\frac{|M|}{\kappa_n}\right)^{1/n},
\]
with precisely the stated homothety equality condition.  The same
derivative identity with $M=K$, using
$|K+tK|=(1+t)^n\kappa_n$, proves
\eqref{eq:eta-normalization}; this identity does not require $K$
to be symmetric.

Now let $y\ne0$ and take $M=[-y,y]$.  Every nonempty fiber of
$K$ parallel to $y$ is a compact interval.  Adding $t[-y,y]$
increases its length by $2t|y|$, without changing its projection
onto $y^\perp$.  Fubini's theorem therefore gives
\[
 |K+t[-y,y]|=|K|+
     2t|y|\operatorname{vol}_{n-1}
                   (\operatorname{proj}_{y^\perp}K)
 =|K|+2t h_{\Pi K}(y).
\]
Since the support function of the segment is $x\mapsto|x\cdot y|$,
the derivative identity yields \eqref{eq:projection-measure}.
The formula is also true at $y=0$.  The projection of a
full-dimensional convex body has positive $(n-1)$-volume in every
direction.  If $\supp\eta_K$ failed to span $\R^n$, a nonzero
vector perpendicular to its span would make the integral in
\eqref{eq:projection-measure} vanish.  This contradiction proves
the spanning property and hence admissibility.

Finally, for $K=B_2^n$ one has $g_K(u)=1$ and
$\nabla g_K(u)=u$ on the sphere, so $\eta_K=\sigma$.
Its projection volume in every unit direction is $\kappa_{n-1}$.
Substituting this into \eqref{eq:projection-measure} proves
\eqref{eq:ball-constant}.
\end{proof}
